\section{Conclusion}

Hanzo Proof of AI makes a stronger claim than AI-themed proof of work. The mined
object is a real, pre-admitted AI workload whose consumer, model, input,
execution policy, proof policy, reward, and payout destination are fixed before
the result. A CPU or GPU provider performs that workload and commits its exact
integer output, transcript, and metering. A-Chain PoAI consensus validates that
the bound useful work was actually done, consumes the global nullifier, and
finalizes one reward receipt.

Ordinary CPUs can validate supported work without repeating the full accelerator
execution: small jobs use exact replay, while challenged matrix products use a
quadratic Freivalds check against a Merkle-bound transcript. This local check is
not the whole security story. Graph binding, coverage, data availability,
verifiable metering, challenge finality, and anti-self-dealing admission are
consensus-critical. Confidential jobs may add TEE evidence; audited P3Q proofs
may eventually replace the optimistic path and settle finalized A-Chain receipt
transitions through Z-Chain.

The economic proposition is direct. Existing AI clouds keep the revenue from
workloads they already serve and may earn additional AI for making those
executions independently accountable. Decentralized providers use the same
interface for open jobs and idle capacity. The network rewards consumed AI
results rather than chain-derived matrix theater.

AI has a hard two-trillion nominal cap: one trillion permanently allocated to
the Hanzo DAO and one trillion reserved for PoAI miners and validators. The
proof-earned pool emits geometrically---500 billion, 250 billion, 125 billion,
and so on---while fee burns can make net supply deflationary as issuance tends
to zero.

The omnichain vision has one work authority and one finalized settlement root.
PoAI consensus exists only on Lux A-Chain. Z-Chain compresses and settles those
finalized receipt transitions through P3Q; it does not originate mining claims.
Lux Quasar and the activated post-quantum validator policy finalize and
authenticate the A/Z settlement root. Any supported chain may consume the
destination-bound receipt exactly once, but no destination may accept raw work
or create a second AI supply. Hanzo is therefore an open useful-compute market
rooted in Lux post-quantum finality, not another isolated AI chain.

The remaining milestone is operational, not rhetorical: a production Zen model
must emit the full deterministic transcript, CPU/GPU vectors must agree, fraud
and withholding must be caught, metering and anti-farming controls must hold,
the two-trillion monetary policy must be enforced, and a P3Q-settled,
PQ-authenticated A-Chain receipt must settle exactly once in multiple
destinations. Mining should remain disabled until that end-to-end test is green.

\paragraph{Cross-references.}
The arithmetic and settlement protocol is specified in the Lux AIVM paper
(\texttt{lux/papers/lux-aivm}) and LP-5000/LP-5200/LP-5300/LP-5301. The intended
post-quantum proof substrate is \texttt{luxfi/p3q}; threshold-signature research
lives in \texttt{luxfi/pulsar} and \texttt{luxfi/corona}. None of those names is
a substitute for an audited, activated end-to-end deployment.

\begin{thebibliography}{9}

\bibitem{freivalds}
R. Freivalds.
\textit{Probabilistic machines can use less running time.}
Information Processing, 1977.

\bibitem{fips204}
National Institute of Standards and Technology.
\textit{FIPS 204: Module-Lattice-Based Digital Signature Standard.}
2024.

\bibitem{matmulpow}
I. Komargodski and O. Weinstein.
\textit{Proofs of Useful Work from Arbitrary Matrix Multiplication.}
arXiv:2504.09971, 2025.

\bibitem{luxpoai}
Z. Kelling.
\textit{AIVM: Proof of AI---The Settlement Primitive for Verifiable Machine
Computation on Lux.}
Lux Papers, 2026.

\end{thebibliography}
